\documentclass[a4paper,11pt]{article}

% --- Packages ---
\usepackage[utf8]{inputenc}
\usepackage[T1]{fontenc}
\usepackage{geometry}
\usepackage{graphicx}
\usepackage{svg}
\usepackage{booktabs}
\usepackage{tabularx}
\usepackage{siunitx}
\usepackage{amsmath}
\usepackage{amssymb}
\usepackage{hyperref}
\usepackage{caption}
\usepackage{fancyhdr}
\usepackage{subcaption}
\usepackage{float}

% --- Configuration ---
\geometry{top=2.5cm, bottom=2.5cm, left=2.5cm, right=2.5cm}

\sisetup{
    output-decimal-marker = {,},
    group-digits = true,
    per-mode = symbol
}

\setlength{\parskip}{8pt}

% SVG Settings
% DO NOT use \svgsetup{convert=true} with the standard 'svg' package.
% Just ensure Inkscape is in PATH and compile with --shell-escape.
% The 'svg' package will automatically call Inkscape to convert .svg to .pdf.

\hypersetup{
    colorlinks=true,
    linkcolor=blue,
    filecolor=magenta,      
    urlcolor=cyan,
    pdftitle={"Morris" 5-DOF Robotic Arm},
    pdfauthor={Benjamin McEwan}
}

% --- DOCUMENT -------------------------------------------------------------------------------------
\begin{document}

\title{\textbf{"Morris" 5-DOF Robotic Arm}}
\author{Benjamin McEwan}
\date{\today}
\maketitle

\tableofcontents


% --- 1: INTRODUCTION ------------------------------------------------------------------------------
\newpage
\section{Introduction}
% \begin{figure}[h!]
%     \centering
%     \includesvg[width=0.5\linewidth]{svg/arm.svg} 
%     \caption{Full arm assembly}
%     \label{fig:arm}
% \end{figure}

Morris 

% --- 2: DRAWINGS ----------------------------------------------------------------------------------
% \newpage
% \section{Drawings}
% \begin{figure}[H]
%     \centering
%     \includesvg[width=0.6\linewidth]{svg/arm_exp.svg}
%     \caption{Exploded arm assembly}
%     \label{fig:assembly}
% \end{figure}

///

% --- 3: BILL OF MATERIALS -------------------------------------------------------------------------
\newpage
\section{Bill of Materials (BOM)}
\begin{table}[H]
    \centering
    \caption{Purchased Parts}
    \label{tab:bom_pp}
    \begin{tabularx}{\textwidth}{@{} l X r @{}}
        \toprule
        \textbf{Description} & \textbf{Qty} \\
        \midrule
        MG996R Servo                                              & 5 \\
        ESP32 Devkit Type-C                                       & 1 \\
        PCA9685 16-Channel 12-Bit PWM Driver                      & 1 \\
        DC5525 to Screw Terminal Connector DC Power Jack (Female) & 1 \\
        10A Adjustable DC-DC Step Down Converter                  & 1 \\
        \bottomrule
    \end{tabularx}
\end{table}

% \begin{table}[H]
%     \centering
%     \caption{Manufactured Parts}
%     \label{tab:bom_manufactured}
%     \begin{tabularx}{\textwidth}{@{} l X r @{}}
%         \toprule
%         \textbf{Part} & \textbf{Material} & \textbf{Qty} \\
%         \midrule
%         \texttt{base}                   & PLA - White & 1 \\
%         \bottomrule
%     \end{tabularx}
% \end{table}

% \begin{table}[H]
%     \centering
%     \caption{Standard Parts}
%     \label{tab:bom_fasteners}
%     \begin{tabularx}{\textwidth}{@{} l X r @{}}
%         \toprule
%         \textbf{Description} & \textbf{Qty} \\
%         \midrule
%         M3x6 Pan Head Screw  & 9  \\
%         M3x10 Pan Head Screw & 22 \\
%         M3x14 Pan Head Screw & 1  \\
%         M3 Nut               & 3  \\
%         \bottomrule
%     \end{tabularx}
% \end{table}


% --- 4: ANALYSIS ----------------------------------------------------------------------------------
\newpage
\section{Analysis}
\subsection{Inverse Kinematics}
When calculating the inverse kinematics (IK) of the Morris robot we will consider the robot's spatial frame $\{\mathrm{s}\}$ to be placed at the center of the base of the robot and aligned such that the $-\mathrm{y}$ direction faces forward, the $-\mathrm{x}$ direction faces right, and the $+\mathrm{z}$ direction faces upwards. 

The $n$ links of the robot are taken as $L_n$ with joints $J_n$. As link 2 has an L-shaped structure it is more specifically denoted by the horizontal and vertical offsets between the joints $J_{2}$ and $J_{3}$, notated as $L_{2x}$ and ${L_2y}$

The following will cover the IK required to position the robot's end-effector at a point $\vec{p} = (p_{x},p_{y},p_{z})$ in 3D space and an orientation $\phi$ in the 2D space of the robot's wrist. The robot is kinematically decoupled into three independent systems. The base pivot $J_{1}$, the main arm $J_{2}$ and $J_{3}$, and the wrist $J_{4}$.

$J_{1}$ is considered a 1R mechanism that acts in the $\mathrm{x}\mathrm{-}\mathrm{y}$ plane that only needs to point towards the target point $\vec{p}$, as such the IK angle for $J_{1}$ is found via
\begin{align*}
    \theta_{1} = \mathrm{atan2}(p_{y},p_{x})
\end{align*}

$J_{2}$, $J_{3}$, and $J_{4}$ are all considered coplanar. In order to apply IK to these joints, the point $\vec{p}$ must be converted into a point $\vec{p'} = (p'_{x}, p'_{y})$ this plane. The $x$ coordinate of this point can be found as magnitude of the $\mathrm{x}\mathrm{-}\mathrm{y}$ plane projection of $\vec{p}$; whilst the $y$ coordinate is simply the $z$ coordinate of $\vec{p}$, therefore
\begin{align*}
    p'_{x} &= \sqrt{ p_{x}^2 + p_{y}^2 } \\
    p'_{y} &= p_{z}
\end{align*}
$J_{2}$ and $J_{3}$ form a 2R mechanism that is decoupled from $J_{4}$. The IK equations of this 2R mechanism must therefore be calculated for a point $\vec{P} = (P_{x}, P_{y})$ which is offset from $\vec{p'}$ by the length $L_{4}$ at the angle $\phi$, the coordinates of this point are found as
\begin{align*}
    P_{x} &= p'_{x} - L_{4}\cos(\phi) \\
    P_{y} &= p'_{y} - L_{4}\sin(\phi) - h
\end{align*}
where $h$ is the vertical offset of $J_{1}$ from the origin of $\{\mathrm{s}\}$.

As link 2 posssesses an L-shaped structure, an "offset angle" $\theta_{0}$ is established, and the length of the effective link $L_{2}$ is calculated from its components,
\begin{align*}
    \theta_0 &= \mathrm{atan2}(L_{2y}, L_{2x}) \\
    L_{2} &= \sqrt{ L_{2x}^2 + L_{2y}^2 }
\end{align*}



\end{document}